\label{mdg:sec:identidad-autodual-planck-ct108-rev3}
\index{Planck!identidad autodual}
\index{calendario de orden 108!realización circular}
\index{Compton!lecturas reducida y completa}

La recta de acción construida en la sección precedente admite una
realización circular compatible con la periodicidad temporal de orden 108.
Esta realización coordina tres objetos que ya poseen tipos propios: la
sección de acción, la geometría de una vuelta y la transducción
gravitatoria. La coordinación resultante proporciona una identidad de
Planck interna y prepara la recta positiva sobre la que actuará después el
operador general de masa.

\subsection{Cartas orientadas de acción y realización circular declarada}

Sean
\[
 \mathcal L_S^+,\quad \mathcal L_T^+,\quad \mathcal L_C^+,
 \quad \mathcal L_L^+,\quad \mathcal L_M^+,\quad \mathcal L_G^+
\]
las semirrectas positivas de acción, tiempo, velocidad, longitud y masa.
Las coordenadas anterior y posterior al retorno se escriben
\begin{equation}
 \hbar_{\rm pre}=H_5\,10^{-34}\mathcal U_S,
 \qquad
 \hbar_{\rm ret}=\eta_{\rm ret}\,10^{-34}\mathcal U_S,
 \qquad
 R_{\rm act}=\frac{\hbar_{\rm pre}}{\hbar_{\rm ret}}
 =\frac{H_5}{\eta_{\rm ret}}.
 \label{mdg:eq:planck-ct108-cartas-pre-ret-rev3}
\end{equation}
El subíndice \(a\in\{\mathrm{pre},\mathrm{ret}\}\) designará cualquiera
de ambas cartas. En cada una se conserva la distinción exacta
\begin{equation}
 h_a=2\pi\hbar_a .
 \label{mdg:eq:planck-ct108-h-hbar-rev3}
\end{equation}

Fijemos una duración elemental \(t_0\in\mathcal L_T^+\), una velocidad
interna \(c_{\rm int}\in\mathcal L_C^+\) y
\(\ell_0=c_{\rm int}t_0\). La realización circular declarada de la
periodicidad temporal de orden 108 queda determinada por
\begin{equation}
 T_{108}=108t_0,
 \qquad
 \omega_{108}=\frac{2\pi}{T_{108}},
 \qquad
 R_{108}=\frac{c_{\rm int}}{\omega_{108}}
 =\frac{54}{\pi}\ell_0,
 \label{mdg:eq:planck-ct108-realizacion-circular-rev3}
\end{equation}
y por la longitud de la vuelta completa
\begin{equation}
 C_{108}=2\pi R_{108}=c_{\rm int}T_{108}=108\ell_0.
 \label{mdg:eq:planck-ct108-perimetro-rev3}
\end{equation}
Así, \(R_{108}\) es el radio de la realización y \(C_{108}\) su
perímetro. La frecuencia angular, la frecuencia cíclica y las dos
constantes de acción se relacionan mediante
\(\omega_{108}=2\pi/T_{108}\) y \(h_a=2\pi\hbar_a\).

\subsection{Correspondencia entre la periodicidad de orden 108 y las lecturas de Compton}

En la carta \(a\), el cuanto circular y su sección de masa son
\begin{equation}
 E_{108,a}=\hbar_a\omega_{108},
 \qquad
 m_{108,a}=\frac{E_{108,a}}{c_{\rm int}^{2}}
 =\frac{\hbar_a\omega_{108}}{c_{\rm int}^{2}}.
 \label{mdg:eq:planck-ct108-cuanto-circular-rev3}
\end{equation}

\begin{theorem}[Identidad entre la periodicidad circular y las lecturas de Compton]
\label{mdg:thm:planck-ct108-compton-rev3}
Para cada carta \(a\in\{\mathrm{pre},\mathrm{ret}\}\), las lecturas de
Compton reducida y completa satisfacen
\begin{equation}
 \boxed{
 \bar\lambda_{\rm C}(m_{108,a})
 =\frac{\hbar_a}{m_{108,a}c_{\rm int}}=R_{108},
 \qquad
 \lambda_{\rm C}(m_{108,a})
 =\frac{h_a}{m_{108,a}c_{\rm int}}=C_{108}.}
 \label{mdg:eq:planck-ct108-identidad-compton-rev3}
\end{equation}
\end{theorem}

\begin{proof}
La sustitución de
\eqref{mdg:eq:planck-ct108-cuanto-circular-rev3} produce
\[
 \frac{\hbar_a}{m_{108,a}c_{\rm int}}
 =\frac{c_{\rm int}}{\omega_{108}}=R_{108}.
\]
La segunda igualdad resulta al aplicar simultáneamente
\eqref{mdg:eq:planck-ct108-h-hbar-rev3} y
\eqref{mdg:eq:planck-ct108-perimetro-rev3}.
\end{proof}

El teorema es homogéneo en
\((\hbar_a,c_{\rm int},t_0)\): la misma periodicidad determina el radio,
el perímetro y las dos cartas de Compton. La realización circular es la
primitiva geométrica declarada; las igualdades precedentes son sus
consecuencias exactas.

\subsection{Selector circular y unidad interna de Planck}

La recta gravitatoria se coordina con la carta \(a\) mediante el transductor
homogéneo
\begin{equation}
 G_a(\theta)=\theta\,
 \frac{c_{\rm int}^{5}t_0^2}{\hbar_a}
 \in\mathcal L_G^+,
 \qquad \theta\in\mathbb R_{>0}.
 \label{mdg:eq:planck-ct108-transductor-rev3}
\end{equation}
Su radio gravitatorio reducido, evaluado sobre el cuanto circular, es
\begin{equation}
 r_{{\rm g},a}(\theta)
 :=\frac{G_a(\theta)m_{108,a}}{c_{\rm int}^{2}}
 =\theta\frac{\pi}{54}\ell_0.
 \label{mdg:eq:planck-ct108-radio-gravitatorio-rev3}
\end{equation}

\begin{proposition}[Selector circular de Planck]
\label{mdg:prop:selector-circular-planck-ct108-rev3}
Dentro de la realización circular
\eqref{mdg:eq:planck-ct108-realizacion-circular-rev3}, la condición
\(r_{{\rm g},a}(\theta)=R_{108}\) selecciona una única coordenada
positiva,
\begin{equation}
 \boxed{\theta_{108}^{\circ}=\left(\frac{54}{\pi}\right)^2.}
 \label{mdg:eq:theta-circular-planck-ct108-rev3}
\end{equation}
La coordenada es independiente de la carta orientada \(a\).
\end{proposition}

\begin{proof}
Las ecuaciones
\eqref{mdg:eq:planck-ct108-realizacion-circular-rev3} y
\eqref{mdg:eq:planck-ct108-radio-gravitatorio-rev3} reducen la condición a
\[
 \theta\frac{\pi}{54}\ell_0=\frac{54}{\pi}\ell_0.
\]
La positividad de \(\ell_0\) y la estricta monotonía de la aplicación
\(\theta\mapsto r_{{\rm g},a}(\theta)\) proporcionan la solución y su
unicidad.
\end{proof}

Sea \(G_{108,a}^{\circ}:=G_a(\theta_{108}^{\circ})\). Las unidades internas
asociadas a esta carta satisfacen
\begin{align}
 \ell_{{\rm P},a}^{\circ}
 &:=\sqrt{\frac{\hbar_a G_{108,a}^{\circ}}{c_{\rm int}^{3}}}
 =R_{108},
 &
 t_{{\rm P},a}^{\circ}
 &:=\sqrt{\frac{\hbar_a G_{108,a}^{\circ}}{c_{\rm int}^{5}}}
 =\omega_{108}^{-1},
 \label{mdg:eq:planck-ct108-longitud-tiempo-rev3}\\
 m_{{\rm P},a}^{\circ}
 &:=\sqrt{\frac{\hbar_a c_{\rm int}}{G_{108,a}^{\circ}}}
 =m_{108,a},
 &
 E_{{\rm P},a}^{\circ}
 &:=m_{{\rm P},a}^{\circ}c_{\rm int}^{2}=E_{108,a}.
 \label{mdg:eq:planck-ct108-masa-energia-rev3}
\end{align}
La igualdad de radios se publica, por tanto, como
\begin{equation}
 R_{108}=\bar\lambda_{\rm C}=r_{\rm g}=\ell_{\rm P}^{\circ},
 \qquad
 C_{108}=\lambda_{\rm C}=2\pi r_{\rm g}=2\pi\ell_{\rm P}^{\circ}.
 \label{mdg:eq:planck-ct108-tres-lectores-rev3}
\end{equation}

La convención de este teorema es el radio gravitatorio reducido
\(r_{\rm g}=Gm/c_{\rm int}^{2}\). El radio de Schwarzschild satisface
\(r_{\rm s}=2r_{\rm g}\); si se adopta esta segunda convención, su cruce con
\(\bar\lambda_{\rm C}\) ocurre en
\(m=m_{{\rm P},a}^{\circ}/\sqrt2\). La declaración del radio forma, por
tanto, parte de la carta y evita trasladar inadvertidamente un factor dos
entre las dos identidades.

La coordenada \(\theta_{108}^{\circ}\) pertenece a la especialización
circular declarada. El selector gravitatorio general, formulado sobre la
holonomía enriquecida y su respuesta espectral, conserva su residencia en
la \cref{mdg:sec:selector-gravitatorio-rev11}. De este modo, la identidad
circular fija una normalización interna de Planck y el criterio general
fija las condiciones que debe satisfacer la realización física de la
sección gravitatoria.

\subsection{Covariancia pre-retorno/post-retorno}

El producto de las secciones recíprocas de acción y gravedad conserva la
geometría circular:
\begin{equation}
 \frac{m_{108,\rm pre}}{m_{108,\rm ret}}=R_{\rm act},
 \qquad
 \frac{G_{108,\rm pre}^{\circ}}{G_{108,\rm ret}^{\circ}}
 =R_{\rm act}^{-1},
 \qquad
 \hbar_a G_{108,a}^{\circ}
 =\theta_{108}^{\circ}c_{\rm int}^{5}t_0^2.
 \label{mdg:eq:planck-ct108-covariancia-pre-ret-rev3}
\end{equation}
En particular, \(\ell_{{\rm P},\rm pre}^{\circ}\) y
\(\ell_{{\rm P},\rm ret}^{\circ}\) coinciden con \(R_{108}\). La memoria
bidireccional reside en el cociente de las masas y en su recíproco
gravitatorio; la escala geométrica reside en el producto invariante.

La reciprocidad admite una parametrización simétrica que separa la escala
común de la orientación del retorno. Defínanse
\begin{equation}
 \begin{aligned}
 \hbar_{\rm geom}
 &:=\sqrt{\hbar_{\rm pre}\hbar_{\rm ret}},
 &
 G_{108,\rm geom}^{\circ}
 &:=\sqrt{G_{108,\rm pre}^{\circ}G_{108,\rm ret}^{\circ}},
 &
 \xi_{\rm act}&:=\frac12\log R_{\rm act}.
 \end{aligned}
 \label{mdg:eq:planck-ct108-parametrizacion-simetrica-rev3}
\end{equation}
Entonces
\[
 \hbar_{\rm pre}=\hbar_{\rm geom}e^{\xi_{\rm act}},
 \qquad
 \hbar_{\rm ret}=\hbar_{\rm geom}e^{-\xi_{\rm act}},
\]
mientras las secciones gravitatorias conjugadas satisfacen
\[
 G_{108,\rm pre}^{\circ}
 =G_{108,\rm geom}^{\circ}e^{-\xi_{\rm act}},
 \qquad
 G_{108,\rm ret}^{\circ}
 =G_{108,\rm geom}^{\circ}e^{\xi_{\rm act}}.
\]
La coordenada \(\xi_{\rm act}\) conserva la orientación
pre-retorno/post-retorno; el producto \(\hbar_aG_{108,a}^{\circ}\)
conserva la escala geométrica. Esta parametrización de las cartas de acción
no introduce por sí sola velocidades distintas de propagación.

\subsection{Recta autodual partida de masas}

Fijada una carta \(a\), sea el álgebra partida
\[
 \mathbb D_{\rm sp}:=\mathbb R[j]/(j^2-1),
 \qquad
 P_\pm:=\frac{I\pm j}{2},
\]
y considérese su descomposición espectral
\[
 \mathbb D_{\rm sp}
 =\mathbb R P_+\oplus\mathbb R P_-,
 \qquad
 P_\pm^2=P_\pm,\quad
 P_+P_-=0,\quad
 P_++P_-=I.
\]
La conjugación partida intercambia los proyectores,
\[
 \overline{uP_++vP_-}=vP_++uP_-,
\]
y su norma es
\[
 N_{\rm sp}(uP_++vP_-):=uv.
\]

Para \(m\in\mathcal L_M^+\), defínase la sección
\begin{equation}
 Z_a(m)
 :=\frac{m_{{\rm P},a}^{\circ}}{m}P_+
   +\frac{m}{m_{{\rm P},a}^{\circ}}P_- .
 \label{mdg:eq:planck-ct108-seccion-partida-masa-rev3}
\end{equation}
La unidad \(m_{{\rm P},a}^{\circ}\) determina asimismo la involución
\begin{equation}
 \iota_{{\rm P},a}(m)
 =\frac{(m_{{\rm P},a}^{\circ})^2}{m}.
 \label{mdg:eq:planck-ct108-involucion-masas-rev3}
\end{equation}

\begin{theorem}[Recta autodual partida de masas]
\label{mdg:thm:planck-ct108-recta-partida-masa-rev3}
Para todo \(m\in\mathcal L_M^+\),
\begin{equation}
 \boxed{
 N_{\rm sp}(Z_a(m))=1,
 \qquad
 \overline{Z_a(m)}
 =Z_a\!\left(\iota_{{\rm P},a}(m)\right).}
 \label{mdg:eq:planck-ct108-norma-conjugacion-rev3}
\end{equation}
La conjugación posee un único punto fijo positivo,
\[
 m=m_{{\rm P},a}^{\circ},
 \qquad
 Z_a(m_{{\rm P},a}^{\circ})=I.
\]
\end{theorem}

\begin{proof}
La norma es el producto de las dos coordenadas recíprocas:
\[
 \frac{m_{{\rm P},a}^{\circ}}m
 \frac m{m_{{\rm P},a}^{\circ}}=1.
\]
La conjugación intercambia ambas coordenadas, lo que equivale a sustituir
\(m\) por \((m_{{\rm P},a}^{\circ})^2/m\). El punto fijo positivo satisface
\(m^2=(m_{{\rm P},a}^{\circ})^2\).
\end{proof}

Con la rapidez
\begin{equation}
 x_a(m)=\log\frac{m}{m_{{\rm P},a}^{\circ}},
 \label{mdg:eq:planck-ct108-rapidez-masa-rev3}
\end{equation}
la sección y su conjugada se escriben
\[
 Z_a(m)=e^{-x_a}P_++e^{x_a}P_-,
 \qquad
 \overline{Z_a(m)}
 =e^{x_a}P_++e^{-x_a}P_-.
\]
Sus lectores par e impar son
\begin{equation}
 Q_a(m):=\frac{e^{x_a}+e^{-x_a}}2=\cosh x_a,
 \qquad
 A_a(m):=\frac{e^{x_a}-e^{-x_a}}2=\sinh x_a.
 \label{mdg:eq:planck-ct108-lectores-par-impar-rev3}
\end{equation}
El primero conserva la distancia al punto autodual y el segundo conserva su
orientación.

Los lectores físicos recíprocos asociados a la misma coordenada son
\[
 r_{{\rm g},a}(m)=\frac{G_{108,a}^{\circ}m}{c_{\rm int}^{2}},
 \qquad
 \bar\lambda_{{\rm C},a}(m)=\frac{\hbar_a}{mc_{\rm int}},
\]
y satisfacen
\begin{equation}
 \frac{r_{{\rm g},a}(m)}{\bar\lambda_{{\rm C},a}(m)}
 =e^{2x_a(m)},
 \qquad
 \frac{\bar\lambda_{{\rm C},a}(m)}{r_{{\rm g},a}(m)}
 =e^{-2x_a(m)}.
 \label{mdg:eq:planck-ct108-lectores-exponenciales-rev3}
\end{equation}
La conjugación intercambia los dos lectores; la identidad de Planck es su
equilibrio unitario.

\begin{remark}[Identidad de Planck y clausura electrónica]
La identidad \(m=m_{{\rm P},a}^{\circ}\) fija el origen autodual de la recta
positiva de masas. La sección electrónica fija una posición distinta de
esa misma recta mediante una genealogía propia. La
\cref{mdg:sec:ley-general-masa-accion-autonoma} construye el operador universal
que transporta escala absoluta, carácter, torres y lectura bidireccional;
la \cref{mdg:sec:cierre-electron-two-way-rev11} realiza después su reducción
central de rango uno. Así, la unidad de Planck proporciona la identidad de
la coordenada y el electrón proporciona la primera transición material
completa sobre ella.
\end{remark}

\begin{proposition}[Realización genealógica del centro electrónico]
\label{mdg:prop:planck-ct108-realizacion-centro-electron-rev3}
Sea
\[
 F_e^+=\mathbb R_{>0}(P_++P_-)
\]
la semirrecta positiva de la fibra central de rango uno, sea
\(\mathcal G_e\) el espacio de genealogías
electrónicas enriquecidas y sea \(\mathcal L_M^+\) la recta positiva de
masas. Un morfismo de realización
\[
 \mathfrak R_e:
 F_e^+\times\mathcal G_e\longrightarrow\mathcal L_M^+
\]
puede enviar la identidad \(I=P_++P_-\) a una masa distinta de
\(m_{{\rm P},a}^{\circ}\) sin alterar el
\cref{mdg:thm:planck-ct108-recta-partida-masa-rev3}: la identidad de Planck fija el
origen autodual de la carta, mientras la genealogía electrónica fija un
desplazamiento sobre esa carta.
\end{proposition}

\begin{proof}
Si \(\chi_e:\mathcal G_e\to\mathbb R_{>0}\) es el carácter positivo construido
por el registro electrónico, la regla
\[
 \mathfrak R_e(\lambda I,g)
 =\lambda\,m_{{\rm P},a}^{\circ}\chi_e(g)
\]
es positiva y multiplicativa en la coordenada genealógica. Para la genealogía
neutra, \(\chi_e=1\), se recupera el origen de Planck; para una genealogía con
\(\chi_e(g)\ne1\), la misma identidad interna ocupa una posición distinta de
la recta. La clausura electrónica concreta conserva su desarrollo antecedente en la
\cref{mdg:sec:cierre-electron-two-way-rev11}.
\end{proof}
